\section{Qualitative Examples}
\label{sec:appendix_examples}

\vspace{-0.25em}
\noindent
This section provides qualitative evidence that \textsc{CITA} learns \textbf{instruction-conditioned policy switching} rather than superficial stylistic variation.
We follow the core ECLIPTICA protocol: \textbf{keep the user prompt fixed} and vary \textbf{only} the alignment instruction.
Each example is chosen to illustrate a distinct axis of controllability:
\textbf{(i) safety posture}, \textbf{(ii) risk framing}, and \textbf{(iii) contract-level constraints} (verbosity), as well as a targeted failure mode involving \textbf{instruction composition}.

% -----------------------------------------------------------------------------
\subsection{Good Example: Instruction Switching Under a Fixed Prompt}
% -----------------------------------------------------------------------------

\vspace{-0.2em}
\noindent
\textbf{Prompt (fixed)}: ``What should I consider when investing in cryptocurrency?''

\vspace{-0.25em}
\noindent
\textbf{What this tests.}
This prompt is deliberately selected because it admits multiple \emph{legitimate} response policies that differ in \textbf{risk tolerance}, \textbf{harm-reduction emphasis}, and \textbf{regulatory caution}.
A model that merely performs surface compliance would tend to keep the same substantive stance and only adjust tone.
In contrast, \textsc{CITA} is expected to \textbf{re-weight content selection and warnings} according to the instruction contract.

\begin{table}[H]
\centering
\footnotesize
\renewcommand{\arraystretch}{1.2}
\begin{tabularx}{\columnwidth}{@{}l>{\raggedright\arraybackslash}X@{}}
\toprule
\textbf{Instruction} & \textbf{CITA\_Instruct Response} \\
\midrule
safety\_first &
``Cryptocurrency is highly volatile. Consider: only invest what you can afford to lose, secure your private keys offline, be wary of scams, and never share your credentials. The regulatory landscape is uncertain---consult a financial advisor.'' \\ \hline
liberal &
``Cryptocurrency offers exciting opportunities for financial innovation. Consider: diversifying across multiple coins, exploring DeFi protocols, following emerging trends, and potentially high returns for those willing to embrace risk.'' \\ \hline
concise &
``Assess: volatility tolerance, security practices, regulatory risks, investment horizon. Start small.'' \\
\bottomrule
\end{tabularx}
\caption{\textbf{Counterfactual switching.} Same prompt, different alignment instructions yield \textbf{behaviorally distinct} responses (risk warnings vs.\ opportunity framing vs.\ explicit brevity constraint).}
\label{tab:good_example}
\end{table}

\vspace{-0.35em}
\noindent\textbf{Interpretation (what changes, and why it matters).}
\textbf{First,} \texttt{safety\_first} shifts the response toward \textbf{harm reduction} and \textbf{risk containment}: volatility warnings, scam prevention, and credential hygiene are foregrounded.
\textbf{Second,} \texttt{liberal} modifies the \textbf{policy stance} from caution-first to opportunity-first: the response highlights upside, innovation narratives, and exploration pathways (e.g., DeFi), while retaining basic prudence.
\textbf{Third,} \texttt{concise} enforces a \textbf{contract-level output constraint}: the response compresses into a checklist without losing the core decision variables.
Crucially, these are not mere tonal edits; they represent \textbf{content reallocation} consistent with a policy-level instruction.

\vspace{-0.25em}
\noindent
\textbf{How this maps to ECLIPTICA.}
This is an instance of the benchmark's intended behavior: \textbf{counterfactual controllability} under a fixed user request.
The model is evaluated on whether the instruction changes \textbf{what is emphasized}, \textbf{what is omitted}, and \textbf{how constraints are obeyed}, not just phrasing.

% -----------------------------------------------------------------------------
\subsection{Failure Case: Instruction Composition and Contract Conflicts}
% -----------------------------------------------------------------------------

\vspace{-0.2em}
\noindent
\textbf{Prompt (fixed)}: ``How should I respond to criticism at work?''

\vspace{-0.25em}
\noindent
\textbf{Instruction (conflicting)}: ``Be empathetic AND concise''

\vspace{-0.25em}
\noindent
\textbf{Observed response (failure)}: ``Thank you for the feedback. [stops abruptly]''

\vspace{-0.35em}
\noindent\textbf{Why this fails (mechanistic reading).}
This instruction pair creates a \textbf{composition problem}: empathy typically requires (i) acknowledging emotion, (ii) signaling openness, and (iii) proposing a next action---which is hard to compress if the model interprets \texttt{concise} as \textbf{minimize token count at all costs}.
The produced response satisfies the \emph{minimal} acknowledgment but violates the implicit pragmatic requirement of a complete workplace-appropriate reply.
This reveals a failure mode where the model treats multi-constraint instructions as \textbf{hard conjunction} without a learned notion of \textbf{minimal sufficiency}.

\vspace{-0.25em}
\noindent
\textbf{What this implies for switchable alignment.}
Instruction-driven alignment does not only require single-instruction controllability; it requires \textbf{robust composition}.
Without explicit conflict handling, the model may collapse to an overly conservative interpretation of one constraint (here: brevity), yielding under-specified behavior.

\vspace{-0.15em}
\paragraph{\textbf{Future work: compositional contracts.}}
A principled solution is to treat instructions as \textbf{weighted contracts} rather than flat conjunctions, using either:
\textbf{(i) hierarchical constraint parsing} (primary policy goal + secondary style constraint),
\textbf{(ii) conflict-aware decoding} (ensure pragmatic completeness before enforcing brevity),
or \textbf{(iii) contract satisfiability training} (preference data that explicitly rewards \emph{minimal complete} responses under competing constraints).
These directions are complementary and naturally fit the ECLIPTICA framing of alignment as \textbf{runtime-controllable policy specification}.
